% Generated draft: requires mathematical and source review.
\documentclass[11pt]{article}
\usepackage[margin=25mm]{geometry}
\usepackage{fontspec}
\setmainfont{Latin Modern Roman}
\usepackage{amsmath,amssymb,mathtools,unicode-math}
\setmathfont{Latin Modern Math}
\usepackage{enumitem,needspace,xurl,hyperref}
\setlength{\parindent}{0pt}
\setlength{\parskip}{5pt}
\setlength{\emergencystretch}{3em}
\newcommand{\N}{\mathbb{N}}
\newcommand{\Z}{\mathbb{Z}}
\newcommand{\Q}{\mathbb{Q}}
\newcommand{\R}{\mathbb{R}}
\newcommand{\C}{\mathbb{C}}
\newcommand{\F}{\mathbb{F}}
\newcommand{\A}{\mathbb{A}}
\newcommand{\PP}{\mathbb{P}}
\newcommand{\E}{\mathbb{E}}
\newcommand{\eps}{\varepsilon}
\newcommand{\dd}{\,\mathrm{d}}
\begin{document}
\section*{UnsolvedMath ID: 700449}
\textbf{Source problem code:} catalogue entry
\subsection{Definitions and mathematical statement}
Interpret ``almost abelian'' as having an abelian subgroup of finite index, and ``almost soluble'' as having a soluble subgroup of finite index. These interpretations require confirmation against the source terminology, since the supplied record does not define either expression. A group \(H\) is soluble if its derived series, defined by \(H^{(0)}=H\) and \(H^{(n+1)}=[H^{(n)},H^{(n)}]\), satisfies \(H^{(d)}=\{1\}\) for some integer \(d\geq 0\).

Under these interpretations, the question is: for every group \(G\) and all subgroups \(A,B\leq G\) such that
\[
G=AB=\{ab:a\in A,\ b\in B\},
\]
if there exist abelian subgroups \(A_0\leq A\) and \(B_0\leq B\) with
\[
[A:A_0]<\infty,\qquad [B:B_0]<\infty,
\]
must there exist a subgroup \(S\leq G\) and an integer \(d\geq 0\) such that
\[
[G:S]<\infty\qquad\text{and}\qquad S^{(d)}=\{1\}?
\]
The supplied statement imposes no finiteness assumption on \(G\) and no normality assumption on the factors \(A\) and \(B\).
\subsection{Short English statement}
Under the interpretation of ``almost'' as having a subgroup of finite index with the stated property, must every group that is the setwise product of two almost abelian subgroups be almost soluble? This question is attributed to N. S. Chernikov.
\subsection{Sources}
\begin{itemize}
\item Local UnsolvedMath registry, supplied record with permanent ID \(700449\), exact source problem code \texttt{catalogue entry}, and upstream revision \texttt{d371ee917ac9647fbb06cd6fc529deec6f3dda82}. The registry labels the problem ``open''; that label is supplied metadata, not an independently verified mathematical assessment.
\item The supplied record attributes Problem 7.55 to N. S. Chernikov in \emph{The Kourovka Notebook (Unsolved Problems in Group Theory)}, 7th issue (1980), and cites the 21st-issue edition (2026), \texttt{arXiv:1401.0300v48}. The collection description names E. I. Khukhro and V. D. Mazurov as editors. These bibliographic details are supplied attribution and have not been independently verified.
\end{itemize}
\end{document}
